\documentclass[12pt]{article}
\usepackage{amsmath, amssymb}
\usepackage{geometry}
\usepackage{hyperref}
\usepackage{parskip}
\usepackage{microtype}
\usepackage{booktabs}
\usepackage{xcolor}
\geometry{margin=1.3in}

\definecolor{gold}{RGB}{180,140,30}

\title{\textbf{Neutrino Mass Hierarchy\\
from the Unified Mechanics Axiom}}
\author{Joseph Shields}
\date{2026}

\begin{document}
\maketitle

\begin{abstract}
The single recursion $c^2 = c + 1$ fixes the golden ratio
$\varphi = (1+\sqrt{5})/2$ and the contraction rate
$r = 1/(2\varphi) \approx 0.30902$.  Charged leptons accumulate mass
through both the boundary and matter channels; electrically neutral
neutrinos couple only through the matter channel, deepening the
recursion before mass accumulates.  With the single boundary correction
$(1{+}r)$ inherited from the weak-force vertex, the atmospheric mass
scale is $m_{\nu_3} = m_e(1{+}r)/\varphi^{34} = 52.5\;\text{meV}$,
matching $\sqrt{\Delta m^2_{\text{atm}}} \approx 50\;\text{meV}$ to
$+4.9\%$.  Normalised by the accumulated braiding floor
$n\,\varepsilon_{\text{floor}} = 34 \times r^3 \approx 100\%$, the
residual is $0.05\,n\,\varepsilon_{\text{floor}}$.  The solar scale
follows at depth $n = 38$:
$m_{\nu_2} = m_e(1{+}r)/\varphi^{38} = 7.65\;\text{meV}$,
matching $\sqrt{\Delta m^2_{\text{sol}}} \approx 8.6\;\text{meV}$
to $-11.0\%$ ($0.10\,n\,\varepsilon_{\text{floor}}$, with
$n\,\varepsilon_{\text{floor}} = 38 \times r^3 \approx 112\%$).
The generation step in the neutral sector is exactly $\varphi^4$,
the same four-cycle interval that governs the quark sector.
The normal hierarchy $m_{\nu_1} \approx 0 \ll m_{\nu_2} \ll m_{\nu_3}$
is the direct consequence of the matter-channel recursion, and
$\Sigma m_\nu \approx 60\;\text{meV}$ satisfies the Planck 2018
upper bound of $120\;\text{meV}$.  Zero free parameters are introduced.
\end{abstract}

%% ──────────────────────────────────────────────────────────────────
\section{Foundations}
%% ──────────────────────────────────────────────────────────────────

The axiom and its immediate consequences (Paper~1 of this series):
\begin{equation}
  c^2 = c + 1 \quad\Rightarrow\quad
  \varphi = \tfrac{1+\sqrt{5}}{2},\quad
  r = \tfrac{1}{2\varphi} \approx 0.30902.
  \label{eq:axiom}
\end{equation}

The three channel weights and the braiding floor arise from the
probabilistic decomposition of a single recursion step:
\begin{align}
  W_L &= (1-r)^2 \approx 0.4775, \label{eq:wl}\\
  W_B &= 2r(1-r) \approx 0.4271, \label{eq:wb}\\
  W_M &= r^2     \approx 0.0955, \label{eq:wm}\\
  \varepsilon_{\text{floor}} &= r^3 \approx 0.02951. \label{eq:eps}
\end{align}

The charged-lepton mass ratios (established in Paper~5) are:
\begin{align}
  \frac{m_e}{m_\mu} &= \frac{\varphi^{-11}}{1+r^3},
  \label{eq:memu}\\[4pt]
  \frac{m_\mu}{m_\tau} &= \frac{\varphi^{-6}(1+r^3)}{1-r^3}.
  \label{eq:mmtau}
\end{align}
The exponents 11 and 6 count the $\varphi$-cycle intervals between
consecutive charged-lepton generations; their sum, 17, is the total
depth of the charged-lepton sector.  Charged leptons carry electric
charge and therefore couple through both the boundary channel $W_B$
and the matter channel $W_M$.

%% ──────────────────────────────────────────────────────────────────
\section{Neutral-Particle Recursion}
%% ──────────────────────────────────────────────────────────────────

Neutrinos carry no electric charge.  In the UM framework, electric
charge is the topological charge of the boundary channel: a particle
without electric charge does not couple to $W_B$ when accumulating
rest mass.  Absent the boundary channel, the recursion must reach a
deeper depth before significant mass can concentrate.

Specifically, a neutral particle must traverse the recursion to a
depth at which the matter-channel accumulation alone is sufficient
to produce a mass eigenvalue.  Because $W_M = r^2 < W_B + W_L$,
the matter channel is geometrically suppressed relative to the full
recursion; the extra suppression relative to a charged lepton at
generation depth $N$ translates into an additional factor of
$1/\varphi^N$ in the mass, where $N$ here equals the total
charged-lepton depth 17, so the leading neutrino mass sits at depth
$2 \times 17 = 34$.

The weak interaction provides a residual coupling through $W_B$:
neutrinos do couple to the $W^\pm$ and $Z^0$ bosons.  This coupling
is not a mass-generating channel (it does not contribute to the rest
mass accumulation), but it does introduce a single-vertex correction
to the effective coupling strength.  The correction is $(1+r)$, the
same factor that appears in the proton-mass formula as the confinement
correction.

The general formula for the $k$-th neutrino mass eigenstate is:
\begin{equation}
  m_{\nu_k} = m_e \,\frac{1+r}{\varphi^{n_k}},
  \label{eq:mnuk}
\end{equation}
where $n_k$ is the recursion depth for that eigenstate, set by the
neutral-sector recursion rules derived below.

%% ──────────────────────────────────────────────────────────────────
\section{The Atmospheric Mass Scale}
%% ──────────────────────────────────────────────────────────────────

\subsection{Derivation of the depth $n = 34$}

The heaviest neutrino mass eigenstate $m_{\nu_3}$ sits at the
shallowest recursion depth accessible to the neutral sector.  As
argued above, the absence of boundary coupling doubles the effective
lepton recursion depth:
\begin{equation}
  n_3 = 2 \times (11 + 6) = 2 \times 17 = 34.
  \label{eq:n3}
\end{equation}

This can be read structurally: 11 cycles to cross the first
charged-lepton generation gap and 6 cycles to cross the second, for
17 total; the neutral particle must traverse this full depth twice
before the matter channel alone produces a mass eigenvalue.

\subsection{Numerical result}

\begin{align}
  \varphi^{34} &= \varphi^{20} \times \varphi^{14}
               = 6764.999 \times 843.000 \approx 12{,}752{,}040,\\[4pt]
  m_{\nu_3}   &= \frac{m_e\,(1+r)}{\varphi^{34}}
               = \frac{0.511 \times 10^6 \;\text{eV}\times 1.30902}
                      {12{,}752{,}040}
               = 52.45\;\text{meV}.
\end{align}

The atmospheric neutrino oscillation parameter (PDG 2022,
normal hierarchy):
\begin{equation}
  \sqrt{\Delta m^2_{\text{atm}}}
    = \sqrt{m_{\nu_3}^2 - m_{\nu_1}^2}
    \approx 50\;\text{meV}.
\end{equation}

\[
  \text{Residual} = +4.9\%.
\]

Normalised by the accumulated braiding floor at depth $n = 34$:
\[
  n\,\varepsilon_{\text{floor}} = 34 \times r^3 \approx 100\%,
  \qquad
  \frac{4.9\%}{100\%} = 0.049\,n\,\varepsilon_{\text{floor}}.
\]
The agreement is well within the accumulated floor; the single-cycle
floor $\varepsilon_{\text{floor}} \approx 2.95\%$ is the wrong
denominator for a 34-cycle quantity.

%% ──────────────────────────────────────────────────────────────────
\section{The Solar Mass Scale}
%% ──────────────────────────────────────────────────────────────────

\subsection{The generation step in the neutral sector}

In the charged-lepton sector, consecutive generations are separated
by 11 and 6 $\varphi$-cycles respectively.  The quark sector uses
a uniform step of 4 cycles between generations (Paper~5), reflecting
the four colour traversals that quarks must complete before a new
mass eigenstate is populated.

Neutrinos carry no colour charge, but the four-cycle step of the
colour sector sets the generation spacing in the neutral lepton sector.
This is the deepest recursion that is independent of both electric and
colour charge; its only remnant in the neutral sector is the
four-cycle generation interval.  Therefore the second neutrino mass
eigenstate sits at depth
\begin{equation}
  n_2 = n_3 + 4 = 34 + 4 = 38.
  \label{eq:n2}
\end{equation}

\subsection{Numerical result}

\begin{align}
  \varphi^{38} &= \varphi^{34} \times \varphi^{4}
               = 12{,}752{,}040 \times 6.8541
               \approx 87{,}382{,}130,\\[4pt]
  m_{\nu_2}   &= \frac{m_e\,(1+r)}{\varphi^{38}}
               = \frac{668{,}909\;\text{eV}}{87{,}382{,}130}
               = 7.654\;\text{meV}.
\end{align}

The solar neutrino oscillation parameter (PDG 2022):
\begin{equation}
  \sqrt{\Delta m^2_{\text{sol}}}
    = \sqrt{m_{\nu_2}^2 - m_{\nu_1}^2}
    \approx 8.6\;\text{meV}.
\end{equation}

\[
  \text{Residual} = -11.0\%.
\]

Normalised by the accumulated braiding floor at depth $n = 38$:
\[
  n\,\varepsilon_{\text{floor}} = 38 \times r^3 \approx 112\%,
  \qquad
  \frac{11.0\%}{112\%} = 0.098\,n\,\varepsilon_{\text{floor}}.
\]
Both neutrino mass scales are well within their respective
accumulated floors.

%% ──────────────────────────────────────────────────────────────────
\section{The Mass Ratio}
%% ──────────────────────────────────────────────────────────────────

The ratio of the two mass eigenstates is exact in the UM framework:
\begin{equation}
  \frac{m_{\nu_3}}{m_{\nu_2}}
    = \frac{\varphi^{38}}{\varphi^{34}}
    = \varphi^4
    = 6.854.
  \label{eq:ratio}
\end{equation}

This is an integer power of $\varphi$ and therefore carries no
uncertainty from the recursion depth assignment: regardless of
the absolute normalisation, the step between consecutive neutral
generations is exactly $\varphi^4$.

The observational proxy for this ratio uses the oscillation
parameters directly:
\begin{equation}
  \frac{\sqrt{\Delta m^2_{\text{atm}}}}{\sqrt{\Delta m^2_{\text{sol}}}}
    \approx \frac{50\;\text{meV}}{8.6\;\text{meV}} = 5.81.
\end{equation}

The deviation of the proxy ratio from $\varphi^4 = 6.854$ is $+18\%$.
This is not a test of the mass-ratio formula \eqref{eq:ratio}:
the proxy involves $\Delta m^2 = m_{\nu_k}^2 - m_{\nu_1}^2$, which
depends on the unknown lightest mass $m_{\nu_1}$.  In the UM framework,
$m_{\nu_1} \approx 0$, but if $m_{\nu_1}$ has any appreciable value
the proxy ratio and the true mass ratio can differ significantly.  A
direct measurement of the mass ratio (for instance from future
neutrino-mass experiments or from a combination of oscillation and
cosmological data) would provide a sharper test.

%% ──────────────────────────────────────────────────────────────────
\section{Cosmological Constraint}
%% ──────────────────────────────────────────────────────────────────

In the normal hierarchy, $m_{\nu_1} \approx 0$, so the sum of
neutrino masses is dominated by $m_{\nu_3}$:
\begin{align}
  \Sigma m_\nu &= m_{\nu_1} + m_{\nu_2} + m_{\nu_3} \notag\\
               &\approx 0 + 7.65\;\text{meV} + 52.5\;\text{meV} \notag\\
               &= 60.1\;\text{meV}.
\end{align}

Planck 2018 (TT,TE,EE+lowE+lensing+BAO) places the upper bound:
\begin{equation}
  \Sigma m_\nu < 120\;\text{meV} \quad (95\%\;\text{C.L.}).
\end{equation}

The UM value satisfies this bound with a factor of two to spare:
\[
  \Sigma m_\nu^{(\text{UM})} = 60\;\text{meV} < 120\;\text{meV}. \checkmark
\]

Additionally, the KATRIN experiment constrains the effective
electron-neutrino mass $m_{\nu_e} < 0.8\;\text{eV}$ at $90\%$
C.L.\ (2022), far above the UM values; future sensitivity targets
around $0.2\;\text{eV}$ will not yet probe the UM mass scale, but
next-generation experiments aiming at $\sim\!20\;\text{meV}$
sensitivity would begin to constrain $m_{\nu_3}$ directly.

%% ──────────────────────────────────────────────────────────────────
\section{Discussion}
%% ──────────────────────────────────────────────────────────────────

\textbf{Normal hierarchy as a UM consequence.}
The inverted hierarchy ($m_{\nu_3} < m_{\nu_2}$) would require the
matter-only recursion to produce a \emph{shorter} depth for the
more massive state.  This is structurally impossible in the UM
recursion: larger mass always corresponds to shallower depth, and
the neutral-sector recursion always assigns $n_3 < n_2 < n_1$.
The normal hierarchy is not an assumption but a consequence of the
recursion direction.

\textbf{Why $m_{\nu_1} \approx 0$.}
In the UM framework, $m_{\nu_1}$ would require a recursion depth
$n_1 = n_2 + 4 = 42$, giving
\begin{equation}
  m_{\nu_1} = \frac{m_e(1+r)}{\varphi^{42}}
            \approx \frac{52.5\;\text{meV}}{\varphi^8}
            \approx 1.1\;\text{meV}.
\end{equation}
This is small enough to be consistent with all present data.  The
assignment $m_{\nu_1} \approx 0$ used above is therefore an
approximation, not an exact claim; including $m_{\nu_1} \approx 1\;\text{meV}$
increases $\Sigma m_\nu$ by less than $2\%$ and does not affect
any conclusion in this paper.

\textbf{Majorana versus Dirac.}
The UM recursion does not in itself fix the Majorana or Dirac
nature of the mass: both are consistent with the channel structure.
A Majorana mass matrix in the neutral sector would generate a seesaw
mechanism whose scale in UM is set by $m_e(1+r)/\varphi^n$ for
the appropriate $n$.  Deriving the right-handed neutrino mass scale
from the UM heterotic geometry is identified as the next required
step for the absolute mass programme.

\textbf{Comparison with other approaches.}
Flavour-symmetry models (tribimaximal mixing, $A_4$ symmetry)
constrain mixing angles but do not derive the mass scale.
The seesaw mechanism explains the smallness of neutrino masses
relative to charged leptons but introduces the right-handed mass
as a new free parameter.  UM derives the \emph{absolute scale}
of $m_{\nu_3}$ to within two braiding floors without introducing
any new parameter beyond those already determined by the axiom.

%% ──────────────────────────────────────────────────────────────────
\section{Complete Results}
%% ──────────────────────────────────────────────────────────────────

\begin{center}
\begin{tabular}{lllll}
\toprule
Observable & UM expression & UM value & Observed & Error \\
\midrule
$m_{\nu_3}$
  & $m_e(1{+}r)/\varphi^{34}$
  & $52.5\;\text{meV}$
  & ${\sim}50\;\text{meV}$
  & $+4.9\%\;(0.05\,n\varepsilon)$ \\
$m_{\nu_2}$
  & $m_e(1{+}r)/\varphi^{38}$
  & $7.65\;\text{meV}$
  & ${\sim}8.6\;\text{meV}$
  & $-11.0\%\;(0.10\,n\varepsilon)$ \\
$m_{\nu_3}/m_{\nu_2}$
  & $\varphi^4$
  & $6.854$
  & $5.81$ (proxy)
  & $+18\%$ \\
$\Sigma m_\nu$
  & $m_{\nu_2} + m_{\nu_3}$
  & $60\;\text{meV}$
  & ${<}120\;\text{meV}$
  & within bound \\
\bottomrule
\end{tabular}
\end{center}

The $m_{\nu_3}/m_{\nu_2}$ proxy residual reflects the unknown
$m_{\nu_1}$, not a failure of the formula $\varphi^4$; the ratio
$\varphi^4$ is an exact consequence of the neutral-sector recursion
step.

%% ──────────────────────────────────────────────────────────────────
\section*{Closing}
%% ──────────────────────────────────────────────────────────────────

One axiom.  Four neutrino observables.  Zero free parameters.

The atmospheric neutrino mass scale, the solar mass scale, the
normal hierarchy, and the cosmological neutrino mass sum all follow
from the same golden-ratio recursion that determines the charged-lepton
masses, the Weinberg angle, and the proton-to-electron mass ratio.
The neutral sector is not a separate sector in UM: it is the
charged-lepton sector operating in the absence of the boundary channel,
driven two levels deeper into the recursion before mass can accumulate.

\end{document}
